\documentclass[aspectratio=169,11pt]{beamer}

\usetheme{Madrid}
\usecolortheme{default}
\usefonttheme{professionalfonts}
\setbeamertemplate{navigation symbols}{}
\setbeamertemplate{footline}[frame number]

\usepackage{amsmath, amssymb, booktabs}

\title{Assignment, Wages, Platforms, And Pricing Rules}
\subtitle{Algorithmic Rules As Labor-Market Institutions}
\author{Special Topic 7: Labor Market Design, Contracting, and Mechanism Design -- Week 4}
\date{}

\begin{document}

\begin{frame}
  \titlepage
\end{frame}

\begin{frame}{Central question}
\begin{itemize}
    \item How do assignment algorithms, platform rules, and wage-setting systems shape labor allocation and worker welfare?
    \item Week 4 moves from contracts inside firms to rules chosen by labor-market intermediaries.
    \item The labor objects are visibility, assignment, wages, commissions, monitoring, reputation, flexibility, risk, and future access to work.
    \item The research standard is rule $\rightarrow$ labor margin $\rightarrow$ counterfactual design $\rightarrow$ welfare incidence.
\end{itemize}
\end{frame}

\begin{frame}{Intermediary objective and rule framework}
\small
\[
  \rho = (V, A, P, M)
\]
\begin{itemize}
    \item $V$: visibility, search order, ranking, recommendations, invitations.
    \item $A$: matching, routing, assignment, scheduling, queue discipline.
    \item $P$: posted wage, bids, commissions, floors, bonuses, dynamic prices, transparency.
    \item $M$: monitoring, ratings, guarantees, dispute rules, deactivation, evaluation.
\end{itemize}
\[
  \max_{\rho} \; \Pi(\rho)
  = \text{fees} + \text{future thickness} + \text{trust}
  - \text{operational and legal costs}
\]
\begin{itemize}
    \item The platform's objective can diverge from worker surplus or social surplus.
\end{itemize}
\end{frame}

\begin{frame}{Worker welfare under platform rules}
\small
\[
W_i(\rho)=
\mathbb{E}[w_i(\rho)]
- c_i(e_i,a_i)
- \lambda_i \operatorname{Var}(w_i(\rho))
+ \phi_i \text{Flex}_i(\rho)
+ q_i(\rho)
- s_i(\rho)
\]
\begin{itemize}
    \item Earnings are only one part of platform worker welfare.
    \item Rules affect time costs, volatility, flexibility, match quality, search burden, and monitoring pressure.
    \item A design can raise fill rates and still shift risk or bargaining power onto workers.
    \item Welfare claims need the relevant worker margin, not just platform output.
\end{itemize}
\end{frame}

\begin{frame}{Assignment, visibility, and ranking}
\small
\begin{tabular}{p{0.25\linewidth} p{0.32\linewidth} p{0.31\linewidth}}
\toprule
Rule & Observable margin & Interpretation risk \\
\midrule
Search ranking & applications, invitations, hires & attention versus quality \\
Recommendation & employer contact, match rate & platform steering versus fit \\
Public signal & hire rate, wages, future jobs & information versus exposure \\
Guarantee & buyer trust, selected workers & credibility versus favoritism \\
\bottomrule
\end{tabular}
\vspace{0.25cm}
\begin{itemize}
    \item Workers and employers rarely see a neutral market menu.
    \item Good papers separate mechanical visibility, belief updating, and equilibrium response.
\end{itemize}
\end{frame}

\begin{frame}{Wages, prices, commissions, and transparency}
\begin{itemize}
    \item Compensation design includes posted wages, bargaining, bids, commissions, wage floors, surge pricing, bonuses, and pay transparency.
    \item Wage information changes search targeting, applicant composition, employer screening, and accepted pay.
    \item Commissions and dynamic prices can allocate workers across times and tasks, but also change risk incidence.
    \item Pay-rule experiments are powerful because the rule itself is varied, not inferred from equilibrium wages alone.
    \item The key question is who absorbs the incidence: workers, employers, consumers, or the platform.
\end{itemize}
\end{frame}

\begin{frame}{Worker welfare, flexibility, and risk}
\begin{itemize}
    \item Flexibility can be a worker amenity when workers control timing and task choice.
    \item Flexibility can also be a risk-shifting channel when workers absorb waiting time, cancellations, low demand, or algorithmic schedule pressure.
    \item Platform welfare objects include earnings, volatility, unpaid time, autonomy, ratings, deactivation risk, and future access to work.
    \item Average hourly pay is incomplete if task availability, risk, and continuation value change at the same time.
    \item Ride-hailing and online labor markets make these nonwage margins empirically visible.
\end{itemize}
\end{frame}

\begin{frame}{Theory-to-applied translation}
\small
\begin{enumerate}
    \item Rule: visibility, assignment, disclosure, wage floor, commission, guarantee, monitoring, or deactivation.
    \item Labor margin: applications, invitations, hires, wage bids, accepted wages, hours, ratings, retention, or future jobs.
    \item Counterfactual: what would users have seen, paid, earned, or faced under the old rule?
    \item Mechanism: information, attention, sorting, bargaining power, risk allocation, or platform steering.
    \item Welfare object: worker surplus, match quality, employer surplus, platform revenue, fairness, flexibility, or access.
\end{enumerate}
\end{frame}

\begin{frame}{Frontier experimental methods}
\small
\begin{tabular}{p{0.33\linewidth} p{0.53\linewidth}}
\toprule
Design & Best use \\
\midrule
Platform A/B tests & ranking, visibility, recommendations, assignment \\
Information disclosure & ratings, histories, wage info, employer messages \\
Recommendation or guarantee randomization & trust, steering, platform credibility \\
Wage-floor or pay-rule experiments & price incidence, participation, bids, fill rates \\
Policy changes or redesigns & transparency, fees, governance, deactivation \\
Survey, conjoint, lab-in-field & fairness, flexibility, algorithm aversion, worker preferences \\
\bottomrule
\end{tabular}
\end{frame}

\begin{frame}{Research Lab design}
\begin{itemize}
    \item Primary anchor: information and visibility for inexperienced workers in an online labor market.
    \item Challenge anchor: platform steering and guarantees; wage-rule extensions through transparency, compensation history, or minimum pay.
    \item Reproduce: compare treated and control workers in a synthetic platform panel.
    \item Diagnose: separate exposure, beliefs, wage changes, future reputation, risk, and welfare incidence.
    \item Transfer: redesign the empirical architecture for ranking, guarantees, wage floors, dynamic pricing, or governance.
    \item Deliverable: rule, labor margin, counterfactual, identification strategy, welfare object, and main threat.
\end{itemize}
\end{frame}

\begin{frame}{Takeaways}
\begin{itemize}
    \item Platforms choose labor-market rules: visibility, assignment, prices, commissions, monitoring, and governance.
    \item The same rule can reduce frictions and shift surplus or risk.
    \item Applied work should identify the precise rule and the labor margin before making welfare claims.
    \item The frontier is empirical design: experiments, policy changes, platform data, and worker-preference evidence.
\end{itemize}
\end{frame}

\end{document}
